% Financiële algebra. Quinten Massys: een geldwisselaar weegt munten, zijn vrouw kijkt mee. Geld lenen en
% wisselen tegen rente was het werk van de Antwerpse wisselaars.
% Motief: de aflossingstabel van een lening met vaste annuïteit in 10 termijnen aan 15 % per termijn (hoog,
% zodat je het ziet). Elke staaf is één termijnbedrag (3 cm): onderaan het kapitaaldeel in fel, bovenaan het
% rentedeel. Met termijnbedrag 1 is het kapitaaldeel A_k = u^(k-1-n): 3 * 1,15^(k-11) cm, telkens 1,15 keer
% het vorige (hoofdstuk 3). De bovenrand van de band is de restschuld S = 1 - u^-(n - t), van links naar
% rechts afgelost, met t = 10 - (21 - x)/2,1 termijnen.
\def\omslagbestand{omslag-geldwisselaar.jpg}
\def\omslagbron{Quinten Massys, \textit{De geldwisselaar en zijn vrouw}, 1514 (publiek domein)}
\def\omslagrand{10.4 + 1.6*(1 - 1.15^(-(21-\x)/2.1))/(1 - 1.15^(-10))}
\def\omslagmotief{%
  \randlijn
  \foreach \k/\h in {1/0.742,2/0.853,3/0.981,4/1.128,5/1.297,6/1.492,7/1.715,8/1.973,9/2.268,10/2.609}{%
    \pgfmathsetmacro{\plek}{13.85+\k*.62}%
    \fill[white,draw=accent,line width=.7pt] (\plek-.22,8.75) rectangle (\plek+.22,11.75);
    \fill[fel,draw=white,line width=.7pt] (\plek-.22,8.75) rectangle (\plek+.22,8.75+\h);
  }
  \draw[white,line width=.8pt] (14.1,8.75) -- (20.6,8.75);
  \node[omslaglabel,anchor=east] at (20.5,12.45) {$A_{k+1} = u \cdot A_k$};
}
